% Options for packages loaded elsewhere
\PassOptionsToPackage{unicode}{hyperref}
\PassOptionsToPackage{hyphens}{url}
\PassOptionsToPackage{dvipsnames,svgnames,x11names}{xcolor}
\documentclass[
  spanish,
  11pt,
]{extarticle}
\usepackage{xcolor}
\usepackage[left=2.5cm, right=2.5cm, top=2.5cm, bottom=2.5cm]{geometry}
\usepackage{amsmath,amssymb}
\setcounter{secnumdepth}{5}
\usepackage{iftex}
\ifPDFTeX
  \usepackage[T1]{fontenc}
  \usepackage[utf8]{inputenc}
  \usepackage{textcomp} % provide euro and other symbols
\else % if luatex or xetex
  \usepackage{unicode-math} % this also loads fontspec
  \defaultfontfeatures{Scale=MatchLowercase}
  \defaultfontfeatures[\rmfamily]{Ligatures=TeX,Scale=1}
\fi
\usepackage{lmodern}
\ifPDFTeX\else
  % xetex/luatex font selection
\fi
% Use upquote if available, for straight quotes in verbatim environments
\IfFileExists{upquote.sty}{\usepackage{upquote}}{}
\IfFileExists{microtype.sty}{% use microtype if available
  \usepackage[]{microtype}
  \UseMicrotypeSet[protrusion]{basicmath} % disable protrusion for tt fonts
}{}
\makeatletter
\@ifundefined{KOMAClassName}{% if non-KOMA class
  \IfFileExists{parskip.sty}{%
    \usepackage{parskip}
  }{% else
    \setlength{\parindent}{0pt}
    \setlength{\parskip}{6pt plus 2pt minus 1pt}}
}{% if KOMA class
  \KOMAoptions{parskip=half}}
\makeatother
\usepackage{longtable,booktabs,array}
\usepackage{calc} % for calculating minipage widths
% Correct order of tables after \paragraph or \subparagraph
\usepackage{etoolbox}
\makeatletter
\patchcmd\longtable{\par}{\if@noskipsec\mbox{}\fi\par}{}{}
\makeatother
% Allow footnotes in longtable head/foot
\IfFileExists{footnotehyper.sty}{\usepackage{footnotehyper}}{\usepackage{footnote}}
\makesavenoteenv{longtable}
\usepackage{graphicx}
\makeatletter
\newsavebox\pandoc@box
\newcommand*\pandocbounded[1]{% scales image to fit in text height/width
  \sbox\pandoc@box{#1}%
  \Gscale@div\@tempa{\textheight}{\dimexpr\ht\pandoc@box+\dp\pandoc@box\relax}%
  \Gscale@div\@tempb{\linewidth}{\wd\pandoc@box}%
  \ifdim\@tempb\p@<\@tempa\p@\let\@tempa\@tempb\fi% select the smaller of both
  \ifdim\@tempa\p@<\p@\scalebox{\@tempa}{\usebox\pandoc@box}%
  \else\usebox{\pandoc@box}%
  \fi%
}
% Set default figure placement to htbp
\def\fps@figure{htbp}
\makeatother
\ifLuaTeX
\usepackage[bidi=basic]{babel}
\else
\usepackage[bidi=default]{babel}
\fi
% get rid of language-specific shorthands (see #6817):
\let\LanguageShortHands\languageshorthands
\def\languageshorthands#1{}
\setlength{\emergencystretch}{3em} % prevent overfull lines
\providecommand{\tightlist}{%
  \setlength{\itemsep}{0pt}\setlength{\parskip}{0pt}}
\usepackage{graphicx}
\usepackage{float}
\usepackage{amssymb}
\usepackage{caption}
\usepackage{esint}
\usepackage{setspace}
\captionsetup[figure]{name=Figura,labelsep=period}
\captionsetup[table]{name=Tabla,labelsep=period}
\setlength{\parskip}{0.5em}
\usepackage{array}
\usepackage{ragged2e}
\usepackage{longtable}
\usepackage{booktabs}
\newcolumntype{L}[1]{>{\RaggedRight\arraybackslash}p{#1}}
\usepackage{etoolbox}
\AtBeginEnvironment{longtable}{\footnotesize}
\AtBeginEnvironment{table}{\footnotesize}
\usepackage{bookmark}
\IfFileExists{xurl.sty}{\usepackage{xurl}}{} % add URL line breaks if available
\urlstyle{same}
\hypersetup{
  pdflang={es},
  colorlinks=true,
  linkcolor={black},
  filecolor={Maroon},
  citecolor={Blue},
  urlcolor={blue},
  pdfcreator={LaTeX via pandoc}}

\author{}
\date{\vspace{-2.5em}}

\begin{document}

\begin{titlepage}
        \centering
        
        \vspace{1cm}
        
        
        \includegraphics[width=\textwidth]{imagenes/tecnologico-de-monterrey-blue.png}
        
        \vspace{2cm}
        
        {\LARGE\bfseries Etapa 3: Análisis Inferencial\par}
        
        \vspace{3cm}
        
        {\LARGE\bfseries Modelación estadística para la toma de decisiones \par}
        
        {\Large Fernanda Meléndez Banda A00843603\par}
        {\Large Carlos Emilio Escobedo Castillo A00843640\par}
        {\Large Eduardo David Vargas Herrera A00843714\par}
        {\Large Rodrigo Fuentes Quintana A00843675\par}
        
        \vspace{0.5cm}
        
        \vfill
        
        {\large 2026-10-15\par}
        
\end{titlepage}

\section{Objetivo del análisis y enfoque
inferencial}\label{objetivo-del-anuxe1lisis-y-enfoque-inferencial}

El objetivo principal del estudio no cambió respecto al análisis
descriptivo, ya que se mantuvo el propósito de describir y comparar la
práctica de lectura de revistas (\(P10\), \(P26\) y \(P11\)) entre
quienes acostumbraban leer materiales considerados por la MOLEC
(\(H\_LEC = 1\)) y quienes no acostumbraban leer pero habían leído en
los últimos meses dichos materiales (\(H\_LEC = 3\)), segmentados por
sexo y condición de actividad (mujeres dedicadas al hogar {[}MH{]},
mujeres que trabajan {[}MT{]} y hombres que trabajan {[}HT{]}). Sin
embargo, el enfoque cambió a una prueba estadística con datos de la
MOLEC 2022-2024. Para lograrlo, se combinó el tipo de ocupación con si
la persona se consideraba lectora (lector/a {[}L{]} frente a no se
percibía {[}NL{]}), formando grupos (como MH-NL o MH-L) que se evaluaron
con los métodos estadísticos explicados a continuación.

\section{Técnicas estadísticas}\label{tuxe9cnicas-estaduxedsticas}

Para evaluar rigurosamente las diferencias entre grupos y celdas, se
emplearon las siguientes herramientas de estadística inferencial:

\subsection{Intervalos de confianza para la diferencia de
medias}\label{intervalos-de-confianza-para-la-diferencia-de-medias}

Se utilizaron para estimar la diferencia entre los promedios de lectores
y no lectores en cada grupo, evaluando el tiempo de lectura (\(P26\)) y
el número de revistas leídas (\(P10\)). Al tratarse de muestras grandes
(\(n, m \ge 40\)) con varianzas poblacionales desconocidas, se aplicó la
distribución normal estándar (\(Z\)) con un nivel de confianza del 96\%:

\[Z = \frac{(\bar{X} - \bar{Y}) - (\mu_X - \mu_Y)}{\sqrt{S_X^2/n + S_Y^2/m}} \approx \mathcal{N}(0, 1)\]

\subsection{Análisis de varianza (ANOVA de dos factores con
interacción)}\label{anuxe1lisis-de-varianza-anova-de-dos-factores-con-interacciuxf3n}

Permitió evaluar simultáneamente si el tipo de ocupación, la
autopercepción lectora y la interacción entre ambos factores generaron
diferencias significativas en las variables cuantitativas (\(P10\) y
\(P26\)). La variabilidad total se desglosó y las medias se contrastaron
mediante el estadístico \(F\) de Snedecor.

\subsection{\texorpdfstring{Prueba de independencia de Chi-cuadrada
(\(\chi^2\))}{Prueba de independencia de Chi-cuadrada (\textbackslash chi\^{}2)}}\label{prueba-de-independencia-de-chi-cuadrada-chi2}

Evaluó la existencia de asociación significativa entre variables
cualitativas, cruzando la autopercepción lectora y el grupo ocupacional
en relación con la lectura (\(P11\)). La técnica contrastó las
frecuencias observadas (\(O_{ij}\)) frente a las esperadas (\(E_{ij}\))
a través del estadístico:

\[Q = \sum_{i=1}^{r} \sum_{j=1}^{c} \frac{(O_{ij} - E_{ij})^2}{E_{ij}} \approx \chi^2_{(r-1)(c-1)}\]

\subsection{Verificación de supuestos del
modelo}\label{verificaciuxf3n-de-supuestos-del-modelo}

Se validaron los requisitos teóricos del ANOVA a través del análisis de
los residuos. La normalidad se comprobó mediante gráficos de cuantiles
teóricos (QQ-plots) para \(P10\) y \(P26\), mientras que la homogeneidad
de varianza y la independencia se evaluaron mediante diagramas de
dispersión de los propios residuos.

\section{Resultados}\label{resultados}

\subsection{Comparaciones evaluadas}\label{comparaciones-evaluadas}

\begin{longtable}[]{@{}
  >{\raggedright\arraybackslash}p{(\linewidth - 6\tabcolsep) * \real{0.0460}}
  >{\raggedright\arraybackslash}p{(\linewidth - 6\tabcolsep) * \real{0.1724}}
  >{\raggedright\arraybackslash}p{(\linewidth - 6\tabcolsep) * \real{0.6207}}
  >{\raggedright\arraybackslash}p{(\linewidth - 6\tabcolsep) * \real{0.1609}}@{}}
\caption{Comparaciones evaluadas}\tabularnewline
\toprule\noalign{}
\begin{minipage}[b]{\linewidth}\raggedright
IC
\end{minipage} & \begin{minipage}[b]{\linewidth}\raggedright
Variable
\end{minipage} & \begin{minipage}[b]{\linewidth}\raggedright
Comparación (celda a - celda b)
\end{minipage} & \begin{minipage}[b]{\linewidth}\raggedright
Claves
\end{minipage} \\
\midrule\noalign{}
\endfirsthead
\toprule\noalign{}
\begin{minipage}[b]{\linewidth}\raggedright
IC
\end{minipage} & \begin{minipage}[b]{\linewidth}\raggedright
Variable
\end{minipage} & \begin{minipage}[b]{\linewidth}\raggedright
Comparación (celda a - celda b)
\end{minipage} & \begin{minipage}[b]{\linewidth}\raggedright
Claves
\end{minipage} \\
\midrule\noalign{}
\endhead
\bottomrule\noalign{}
\endlastfoot
IC1 & P26 (min) & Mujeres del hogar: no lectoras - lectoras & MH-NL -
MH-L \\
IC2 & P26 (min) & Mujeres que trabajan: no lectoras - lectoras & MT-NL -
MT-L \\
IC3 & P26 (min) & Hombres que trabajan: no lectores - lectores & HT-NL -
HT-L \\
IC4 & P26 (min) & No lectoras: mujeres del hogar - mujeres que trabajan
& MH-NL - MT-NL \\
IC5 & P26 (min) & Lectores: hombres que trabajan - mujeres que trabajan
& HT-L - MT-L \\
IC6 & P10 (revistas) & Mujeres del hogar: no lectoras - lectoras & MH-NL
- MH-L \\
IC7 & P10 (revistas) & Mujeres que trabajan: no lectoras - lectoras &
MT-NL - MT-L \\
IC8 & P10 (revistas) & Hombres que trabajan: no lectores - lectores &
HT-NL - HT-L \\
\end{longtable}

\subsection{Intervalos de confianza para la diferencia de
medias}\label{intervalos-de-confianza-para-la-diferencia-de-medias-1}

\begin{longtable}[]{@{}
  >{\raggedright\arraybackslash}p{(\linewidth - 18\tabcolsep) * \real{0.0519}}
  >{\raggedright\arraybackslash}p{(\linewidth - 18\tabcolsep) * \real{0.1948}}
  >{\raggedleft\arraybackslash}p{(\linewidth - 18\tabcolsep) * \real{0.0519}}
  >{\raggedleft\arraybackslash}p{(\linewidth - 18\tabcolsep) * \real{0.0519}}
  >{\raggedleft\arraybackslash}p{(\linewidth - 18\tabcolsep) * \real{0.1039}}
  >{\raggedleft\arraybackslash}p{(\linewidth - 18\tabcolsep) * \real{0.1039}}
  >{\raggedleft\arraybackslash}p{(\linewidth - 18\tabcolsep) * \real{0.1039}}
  >{\raggedleft\arraybackslash}p{(\linewidth - 18\tabcolsep) * \real{0.1039}}
  >{\raggedleft\arraybackslash}p{(\linewidth - 18\tabcolsep) * \real{0.1039}}
  >{\raggedright\arraybackslash}p{(\linewidth - 18\tabcolsep) * \real{0.1299}}@{}}
\caption{Intervalos de confianza al 96 \% para la diferencia de medias
(celda a - celda b)}\tabularnewline
\toprule\noalign{}
\begin{minipage}[b]{\linewidth}\raggedright
IC
\end{minipage} & \begin{minipage}[b]{\linewidth}\raggedright
Variable
\end{minipage} & \begin{minipage}[b]{\linewidth}\raggedleft
n1
\end{minipage} & \begin{minipage}[b]{\linewidth}\raggedleft
n2
\end{minipage} & \begin{minipage}[b]{\linewidth}\raggedleft
Media\_1
\end{minipage} & \begin{minipage}[b]{\linewidth}\raggedleft
Media\_2
\end{minipage} & \begin{minipage}[b]{\linewidth}\raggedleft
Dif
\end{minipage} & \begin{minipage}[b]{\linewidth}\raggedleft
LI
\end{minipage} & \begin{minipage}[b]{\linewidth}\raggedleft
LS
\end{minipage} & \begin{minipage}[b]{\linewidth}\raggedright
Incluye 0
\end{minipage} \\
\midrule\noalign{}
\endfirsthead
\toprule\noalign{}
\begin{minipage}[b]{\linewidth}\raggedright
IC
\end{minipage} & \begin{minipage}[b]{\linewidth}\raggedright
Variable
\end{minipage} & \begin{minipage}[b]{\linewidth}\raggedleft
n1
\end{minipage} & \begin{minipage}[b]{\linewidth}\raggedleft
n2
\end{minipage} & \begin{minipage}[b]{\linewidth}\raggedleft
Media\_1
\end{minipage} & \begin{minipage}[b]{\linewidth}\raggedleft
Media\_2
\end{minipage} & \begin{minipage}[b]{\linewidth}\raggedleft
Dif
\end{minipage} & \begin{minipage}[b]{\linewidth}\raggedleft
LI
\end{minipage} & \begin{minipage}[b]{\linewidth}\raggedleft
LS
\end{minipage} & \begin{minipage}[b]{\linewidth}\raggedright
Incluye 0
\end{minipage} \\
\midrule\noalign{}
\endhead
\bottomrule\noalign{}
\endlastfoot
IC1 & P26 (min) & 48 & 219 & 23.542 & 35.721 & -12.180 & -16.339 &
-8.021 & No \\
IC2 & P26 (min) & 80 & 335 & 26.500 & 40.910 & -14.410 & -21.082 &
-7.739 & No \\
IC3 & P26 (min) & 58 & 406 & 28.397 & 46.761 & -18.365 & -26.460 &
-10.269 & No \\
IC4 & P26 (min) & 48 & 80 & 23.542 & 26.500 & -2.958 & -9.428 & 3.511 &
Sí \\
IC5 & P26 (min) & 406 & 335 & 46.761 & 40.910 & 5.851 & 0.665 & 11.036 &
No \\
IC6 & P10 (revistas) & 48 & 219 & 2.271 & 3.288 & -1.017 & -1.943 &
-0.090 & No \\
IC7 & P10 (revistas) & 80 & 335 & 2.312 & 3.260 & -0.947 & -1.484 &
-0.411 & No \\
IC8 & P10 (revistas) & 58 & 406 & 3.086 & 3.633 & -0.547 & -1.525 &
0.432 & Sí \\
\end{longtable}

\begin{figure}[H]

{\centering \includegraphics[width=1\linewidth]{Etapa3_Inferencia_files/figure-latex/fig-forest-1} 

}

\caption{Diferencias de medias con IC al 96 \%. En rojo, los intervalos que no incluyen el 0.}\label{fig:fig-forest}
\end{figure}

\subsubsection{Interpretación}\label{interpretaciuxf3n}

\subsection{Efecto del grupo y de la autopercepción: ANOVA de dos
factores}\label{efecto-del-grupo-y-de-la-autopercepciuxf3n-anova-de-dos-factores}

\subsubsection{\texorpdfstring{Número de revistas
(\texttt{P10})}{Número de revistas (P10)}}\label{nuxfamero-de-revistas-p10}

\begin{longtable}[]{@{}lrrrrll@{}}
\caption{ANOVA con interacción: número de revistas (P10)}\tabularnewline
\toprule\noalign{}
Fuente & gl & SC & CM & F & p & Decision \\
\midrule\noalign{}
\endfirsthead
\toprule\noalign{}
Fuente & gl & SC & CM & F & p & Decision \\
\midrule\noalign{}
\endhead
\bottomrule\noalign{}
\endlastfoot
Grupo & 2 & 62.87 & 31.434 & 1.445 & 0.2362 & No rechaza H0 \\
Aut & 1 & 107.49 & 107.491 & 4.941 & 0.0264 & Rechaza H0 \\
Grupo x Aut & 2 & 6.33 & 3.165 & 0.145 & 0.8646 & No rechaza H0 \\
Residuals & 1140 & 24798.84 & 21.753 & & & \\
Total & 1145 & 24975.53 & & & & \\
\end{longtable}

\begin{longtable}[]{@{}lrrrrll@{}}
\caption{Modelo final (solo factores significativos):
P10}\tabularnewline
\toprule\noalign{}
Fuente & gl & SC & CM & F & p & Decision \\
\midrule\noalign{}
\endfirsthead
\toprule\noalign{}
Fuente & gl & SC & CM & F & p & Decision \\
\midrule\noalign{}
\endhead
\bottomrule\noalign{}
\endlastfoot
Aut & 1 & 120.92 & 120.920 & 5.566 & 0.0185 & Rechaza H0 \\
Residuals & 1144 & 24854.60 & 21.726 & & & \\
Total & 1145 & 24975.53 & & & & \\
\end{longtable}

\begin{longtable}[]{@{}lrrr@{}}
\caption{Efecto estimado por autopercepción (media del nivel - media
general): P10}\tabularnewline
\toprule\noalign{}
Nivel & n & Media & Efecto \\
\midrule\noalign{}
\endfirsthead
\toprule\noalign{}
Nivel & n & Media & Efecto \\
\midrule\noalign{}
\endhead
\bottomrule\noalign{}
\endlastfoot
Se percibe lector & 960 & 3.424 & 0.143 \\
No se percibe lector & 186 & 2.543 & -0.738 \\
\end{longtable}

\paragraph{Interpretación}\label{interpretaciuxf3n-1}

\subsubsection{\texorpdfstring{Minutos continuos de lectura
(\texttt{P26})}{Minutos continuos de lectura (P26)}}\label{minutos-continuos-de-lectura-p26}

\begin{longtable}[]{@{}
  >{\raggedright\arraybackslash}p{(\linewidth - 12\tabcolsep) * \real{0.1765}}
  >{\raggedleft\arraybackslash}p{(\linewidth - 12\tabcolsep) * \real{0.0735}}
  >{\raggedleft\arraybackslash}p{(\linewidth - 12\tabcolsep) * \real{0.1618}}
  >{\raggedleft\arraybackslash}p{(\linewidth - 12\tabcolsep) * \real{0.1471}}
  >{\raggedleft\arraybackslash}p{(\linewidth - 12\tabcolsep) * \real{0.1029}}
  >{\raggedright\arraybackslash}p{(\linewidth - 12\tabcolsep) * \real{0.1324}}
  >{\raggedright\arraybackslash}p{(\linewidth - 12\tabcolsep) * \real{0.2059}}@{}}
\caption{ANOVA con interacción: minutos continuos de lectura
(P26)}\tabularnewline
\toprule\noalign{}
\begin{minipage}[b]{\linewidth}\raggedright
Fuente
\end{minipage} & \begin{minipage}[b]{\linewidth}\raggedleft
gl
\end{minipage} & \begin{minipage}[b]{\linewidth}\raggedleft
SC
\end{minipage} & \begin{minipage}[b]{\linewidth}\raggedleft
CM
\end{minipage} & \begin{minipage}[b]{\linewidth}\raggedleft
F
\end{minipage} & \begin{minipage}[b]{\linewidth}\raggedright
p
\end{minipage} & \begin{minipage}[b]{\linewidth}\raggedright
Decision
\end{minipage} \\
\midrule\noalign{}
\endfirsthead
\toprule\noalign{}
\begin{minipage}[b]{\linewidth}\raggedright
Fuente
\end{minipage} & \begin{minipage}[b]{\linewidth}\raggedleft
gl
\end{minipage} & \begin{minipage}[b]{\linewidth}\raggedleft
SC
\end{minipage} & \begin{minipage}[b]{\linewidth}\raggedleft
CM
\end{minipage} & \begin{minipage}[b]{\linewidth}\raggedleft
F
\end{minipage} & \begin{minipage}[b]{\linewidth}\raggedright
p
\end{minipage} & \begin{minipage}[b]{\linewidth}\raggedright
Decision
\end{minipage} \\
\midrule\noalign{}
\endhead
\bottomrule\noalign{}
\endlastfoot
Grupo & 2 & 21709.02 & 10854.509 & 10.978 & \textless{} 0.0001 & Rechaza
H0 \\
Aut & 1 & 35459.65 & 35459.646 & 35.865 & \textless{} 0.0001 & Rechaza
H0 \\
Grupo x Aut & 2 & 907.05 & 453.525 & 0.459 & 0.6322 & No rechaza H0 \\
Residuals & 1140 & 1127124.94 & 988.706 & & & \\
Total & 1145 & 1185200.66 & & & & \\
\end{longtable}

\begin{longtable}[]{@{}lrrrrll@{}}
\caption{Modelo final (solo factores significativos):
P26}\tabularnewline
\toprule\noalign{}
Fuente & gl & SC & CM & F & p & Decision \\
\midrule\noalign{}
\endfirsthead
\toprule\noalign{}
Fuente & gl & SC & CM & F & p & Decision \\
\midrule\noalign{}
\endhead
\bottomrule\noalign{}
\endlastfoot
Grupo & 2 & 21709.02 & 10854.509 & 10.989 & \textless{} 0.0001 & Rechaza
H0 \\
Aut & 1 & 35459.65 & 35459.646 & 35.899 & \textless{} 0.0001 & Rechaza
H0 \\
Residuals & 1142 & 1128031.99 & 987.769 & & & \\
Total & 1145 & 1185200.66 & & & & \\
\end{longtable}

\begin{longtable}[]{@{}lrrr@{}}
\caption{Efecto estimado por grupo ocupacional (media del grupo - media
general): P26}\tabularnewline
\toprule\noalign{}
Nivel & n & Media & Efecto \\
\midrule\noalign{}
\endfirsthead
\toprule\noalign{}
Nivel & n & Media & Efecto \\
\midrule\noalign{}
\endhead
\bottomrule\noalign{}
\endlastfoot
Mujer hogar & 267 & 33.532 & -6.093 \\
Mujer trabaja & 415 & 38.133 & -1.492 \\
Hombre trabaja & 464 & 44.466 & 4.841 \\
\end{longtable}

\begin{longtable}[]{@{}lrrr@{}}
\caption{Efecto estimado por autopercepción (media del nivel - media
general): P26}\tabularnewline
\toprule\noalign{}
Nivel & n & Media & Efecto \\
\midrule\noalign{}
\endfirsthead
\toprule\noalign{}
Nivel & n & Media & Efecto \\
\midrule\noalign{}
\endhead
\bottomrule\noalign{}
\endlastfoot
Se percibe lector & 960 & 42.201 & 2.576 \\
No se percibe lector & 186 & 26.328 & -13.297 \\
\end{longtable}

\begin{longtable}[]{@{}lrrrl@{}}
\caption{Comparaciones de Tukey entre grupos: P26}\tabularnewline
\toprule\noalign{}
Comparacion & Dif & LI & LS & p\_ajustado \\
\midrule\noalign{}
\endfirsthead
\toprule\noalign{}
Comparacion & Dif & LI & LS & p\_ajustado \\
\midrule\noalign{}
\endhead
\bottomrule\noalign{}
\endlastfoot
Mujer trabaja-Mujer hogar & 4.601 & -1.400 & 10.602 & 0.1492 \\
Hombre trabaja-Mujer hogar & 10.934 & 5.058 & 16.809 & \textless{}
0.0001 \\
Hombre trabaja-Mujer trabaja & 6.333 & 1.165 & 11.501 & 0.0082 \\
\end{longtable}

\paragraph{Interpretación}\label{interpretaciuxf3n-2}

\subsubsection{Resumen de los modelos}\label{resumen-de-los-modelos}

\begin{longtable}[]{@{}
  >{\raggedright\arraybackslash}p{(\linewidth - 10\tabcolsep) * \real{0.1724}}
  >{\raggedright\arraybackslash}p{(\linewidth - 10\tabcolsep) * \real{0.1609}}
  >{\raggedright\arraybackslash}p{(\linewidth - 10\tabcolsep) * \real{0.1034}}
  >{\raggedright\arraybackslash}p{(\linewidth - 10\tabcolsep) * \real{0.1954}}
  >{\raggedright\arraybackslash}p{(\linewidth - 10\tabcolsep) * \real{0.1494}}
  >{\raggedright\arraybackslash}p{(\linewidth - 10\tabcolsep) * \real{0.2184}}@{}}
\caption{Resumen de los modelos de dos factores (alpha =
0.04)}\tabularnewline
\toprule\noalign{}
\begin{minipage}[b]{\linewidth}\raggedright
Variable
\end{minipage} & \begin{minipage}[b]{\linewidth}\raggedright
p\_interaccion
\end{minipage} & \begin{minipage}[b]{\linewidth}\raggedright
p\_grupo
\end{minipage} & \begin{minipage}[b]{\linewidth}\raggedright
p\_autopercepcion
\end{minipage} & \begin{minipage}[b]{\linewidth}\raggedright
Modelo\_final
\end{minipage} & \begin{minipage}[b]{\linewidth}\raggedright
Varianza\_explicada
\end{minipage} \\
\midrule\noalign{}
\endfirsthead
\toprule\noalign{}
\begin{minipage}[b]{\linewidth}\raggedright
Variable
\end{minipage} & \begin{minipage}[b]{\linewidth}\raggedright
p\_interaccion
\end{minipage} & \begin{minipage}[b]{\linewidth}\raggedright
p\_grupo
\end{minipage} & \begin{minipage}[b]{\linewidth}\raggedright
p\_autopercepcion
\end{minipage} & \begin{minipage}[b]{\linewidth}\raggedright
Modelo\_final
\end{minipage} & \begin{minipage}[b]{\linewidth}\raggedright
Varianza\_explicada
\end{minipage} \\
\midrule\noalign{}
\endhead
\bottomrule\noalign{}
\endlastfoot
P10 (revistas) & 0.8646 & 0.2357 & 0.0263 & Aut & 0.48 \% \\
P26 (minutos) & 0.6322 & \textless{} 0.0001 & \textless{} 0.0001 &
Grupo, Aut & 4.82 \% \\
\end{longtable}

\subsection{\texorpdfstring{Motivo de lectura (\texttt{P11}):
independencia de dos variables
categóricas}{Motivo de lectura (P11): independencia de dos variables categóricas}}\label{motivo-de-lectura-p11-independencia-de-dos-variables-categuxf3ricas}

\begin{longtable}[]{@{}lrr@{}}
\caption{Frecuencias observadas: autopercepción x motivo de
lectura}\tabularnewline
\toprule\noalign{}
& Por gusto & Otro motivo \\
\midrule\noalign{}
\endfirsthead
\toprule\noalign{}
& Por gusto & Otro motivo \\
\midrule\noalign{}
\endhead
\bottomrule\noalign{}
\endlastfoot
Se percibe lector & 524 & 436 \\
No se percibe lector & 133 & 53 \\
\end{longtable}

\begin{longtable}[]{@{}lrr@{}}
\caption{Frecuencias esperadas bajo independencia}\tabularnewline
\toprule\noalign{}
& Por gusto & Otro motivo \\
\midrule\noalign{}
\endfirsthead
\toprule\noalign{}
& Por gusto & Otro motivo \\
\midrule\noalign{}
\endhead
\bottomrule\noalign{}
\endlastfoot
Se percibe lector & 550.37 & 409.63 \\
No se percibe lector & 106.63 & 79.37 \\
\end{longtable}

\begin{longtable}[]{@{}lll@{}}
\caption{Porcentaje que lee revistas por gusto o entretenimiento, por
grupo y autopercepción}\tabularnewline
\toprule\noalign{}
Grupo & No se percibe lector & Se percibe lector \\
\midrule\noalign{}
\endfirsthead
\toprule\noalign{}
Grupo & No se percibe lector & Se percibe lector \\
\midrule\noalign{}
\endhead
\bottomrule\noalign{}
\endlastfoot
Mujer hogar & 83.3 \% (n = 48) & 67.1 \% (n = 219) \\
Mujer trabaja & 75.0 \% (n = 80) & 52.5 \% (n = 335) \\
Hombre trabaja & 56.9 \% (n = 58) & 49.5 \% (n = 406) \\
\end{longtable}

\begin{longtable}[]{@{}
  >{\raggedright\arraybackslash}p{(\linewidth - 14\tabcolsep) * \real{0.3874}}
  >{\raggedleft\arraybackslash}p{(\linewidth - 14\tabcolsep) * \real{0.0450}}
  >{\raggedleft\arraybackslash}p{(\linewidth - 14\tabcolsep) * \real{0.0270}}
  >{\raggedleft\arraybackslash}p{(\linewidth - 14\tabcolsep) * \real{0.1261}}
  >{\raggedleft\arraybackslash}p{(\linewidth - 14\tabcolsep) * \real{0.0631}}
  >{\raggedright\arraybackslash}p{(\linewidth - 14\tabcolsep) * \real{0.0811}}
  >{\raggedleft\arraybackslash}p{(\linewidth - 14\tabcolsep) * \real{0.1441}}
  >{\raggedright\arraybackslash}p{(\linewidth - 14\tabcolsep) * \real{0.1261}}@{}}
\caption{Pruebas ji cuadrada de independencia con el motivo de lectura
(P11)}\tabularnewline
\toprule\noalign{}
\begin{minipage}[b]{\linewidth}\raggedright
Prueba
\end{minipage} & \begin{minipage}[b]{\linewidth}\raggedleft
n
\end{minipage} & \begin{minipage}[b]{\linewidth}\raggedleft
gl
\end{minipage} & \begin{minipage}[b]{\linewidth}\raggedleft
Valor crítico
\end{minipage} & \begin{minipage}[b]{\linewidth}\raggedleft
q
\end{minipage} & \begin{minipage}[b]{\linewidth}\raggedright
p
\end{minipage} & \begin{minipage}[b]{\linewidth}\raggedleft
Esperada mínima
\end{minipage} & \begin{minipage}[b]{\linewidth}\raggedright
Decisión
\end{minipage} \\
\midrule\noalign{}
\endfirsthead
\toprule\noalign{}
\begin{minipage}[b]{\linewidth}\raggedright
Prueba
\end{minipage} & \begin{minipage}[b]{\linewidth}\raggedleft
n
\end{minipage} & \begin{minipage}[b]{\linewidth}\raggedleft
gl
\end{minipage} & \begin{minipage}[b]{\linewidth}\raggedleft
Valor crítico
\end{minipage} & \begin{minipage}[b]{\linewidth}\raggedleft
q
\end{minipage} & \begin{minipage}[b]{\linewidth}\raggedright
p
\end{minipage} & \begin{minipage}[b]{\linewidth}\raggedleft
Esperada mínima
\end{minipage} & \begin{minipage}[b]{\linewidth}\raggedright
Decisión
\end{minipage} \\
\midrule\noalign{}
\endhead
\bottomrule\noalign{}
\endlastfoot
Autopercepción x motivo (todos los grupos) & 1146 & 1 & 4.218 & 18.239 &
\textless{} 0.0001 & 79.4 & Rechaza H0 \\
Grupo x motivo (todos los casos) & 1146 & 2 & 6.438 & 26.689 &
\textless{} 0.0001 & 113.9 & Rechaza H0 \\
Autopercepción x motivo: Mujer hogar & 267 & 1 & 4.218 & 4.930 & 0.0264
& 14.4 & Rechaza H0 \\
Autopercepción x motivo: Mujer trabaja & 415 & 1 & 4.218 & 13.284 &
0.0003 & 34.5 & Rechaza H0 \\
Autopercepción x motivo: Hombre trabaja & 464 & 1 & 4.218 & 1.108 &
0.2924 & 28.8 & No rechaza H0 \\
\end{longtable}

\subsubsection{Interpretación}\label{interpretaciuxf3n-3}

\subsection{Verificación de
supuestos}\label{verificaciuxf3n-de-supuestos}

\begin{figure}[H]

{\centering \includegraphics[width=1\linewidth]{Etapa3_Inferencia_files/figure-latex/fig-supuestos-1} 

}

\caption{Diagnóstico de los residuos de los modelos finales: normalidad (arriba) y residuos contra valores ajustados (abajo).}\label{fig:fig-supuestos}
\end{figure}

\subsubsection{Interpretación}\label{interpretaciuxf3n-4}

\section{Conclusiones}\label{conclusiones}

\end{document}
